% !TeX root = ../../main.tex
\documentclass[../../main.tex]{subfiles}
\begin{document}

\subsubsection{Atomization~0}

The atomizer has to be prepared in advance for the first trial of gas atomization. The temperature of the melt has to be measured during the process, and the optical thermal sensor kit is used for this purpose. The first task is to set up the sensor with the right driver and software and testing the response in front of en opened furnace.

\paragraph{Tools / method / protocol}
\label{}

\subparagraph{Preparation of the atomizer}
\label{}

2026/03/17

\begin{enumerate}
    \item Preparing the crucible, check there is no crack. In case of cracks do some cement and fix them. Dry in the oven the proper way.
    \item Put the crucible in the melting chamber. Make sure that it sit properly on the spacers.
    \item Put the feedstock rods in the crucible. They should stand vertically letting enough space in the center for the stopper rod. They should be able to extend due to the heat without touching the walls pushing on the crucible inner side (risk of breaking the crucible). This corresponds to approximately 10 kg of stainless steel bars.
    \item Put some insolation to cover the top of the crucible and the feedstock rods. An aperture should be left to let the stopper rod go through and also to let the optical sensor look the melt.
    \todo{$\succ$ Design a cap of solid foam to replace this insolation.}
    \item Put the main cover thanks to the crane and screw it properly to compress the O-ring.
    \item Put the main flange and screw it properly to compress the O-ring.
    \item Put the stopper rod and fasten it. Raise the horizontal support compressing a bit both of the spring and fasten the rod.
    \item Put the optical sensor and fasten it. The sensor should be looking through the dedicated window to the melt.

    \item Check the all the valves.
    \item Start the vacuum pump until reaching 0.074 bar or lower.
    \item Close the vacuum valves, start the N2 flow and switch off the vacuum pump. Pressure should raise to atmospheric pressure.
    \item Start the vacuum pump again until reaching 0.074 bar or lower (one N2 flush is done).
    \item Close the vacuum valves, start the N2 flow and switch off the vacuum pump. Pressure should be just above atmospheric pressure. (N2 pressure should be 0.15 / 0.18 bar)
    \item Oxygen percentage should be 0.2 / 0.3 in the N2
    \todo{$\succ$ How to calibrate these sensors?}

\end{enumerate}

\subparagraph{Powder measurement}
\label{}

\begin{enumerate}
    \item Collect the powder in a heat proof container
    \item Put it in an oven during 12 hours at 120$^\circ$C to dry it.
    \item Separate the flaxes from the power by choosing the right sieve size.
    \item Then use the standard sieves to separate the powder in different size classes and weight them.
    \item Clean the sieves and the container.
\end{enumerate}

\paragraph{Results}

Thanks to the latest improvement of the atomizer, the vacuum rate was very good, reaching 0.074 bar and the oxygen level was around 0.2\% in the N2 after the two flushes, once the atomizer was under N2 atmosphere, ready for atomization. A very small amount of oxidation were floating on the melt surface, which shows that the oxygen level was quite low.\\

Approximately 10 kg of stainless steel bars were put in the crucible as a feedstock. During this atomization1, 3 main problems appeared:
\begin{itemize}
    \item A leakage between the metal nozzle and the crucible allowed the liquid metal to go in between
    \item A freeze-off appeared at the end of the atomization
    \item The cyclone get stuck very early in the atomization, preventing the powder to go in the powder container. Instead, the powder was drained with the gaz in both of the water filters.
\end{itemize}

After the atomization, the weight of the obtained materials have been weighed:
\begin{itemize}
    \item The metal freeze attached to the nozzle: $\approx$ 500g
    \item Flaxes in the container: $\approx$ 3040 g
    \item Powder in the container [125 $<$ x $<$ 63]: 806 g
    \item Powder in the container [x $<$ 63]: 121 g
    \item Flaxes and powder in the cyclone: $\approx$ 400 g
\end{itemize}






\end{document}
